\honest{Adaptation-Executers, not Fitness-Maximizers}{Eliezer Yudkowsky}{2007}

My thesis is a sentence from John Tooby and Leda Cosmides: ``Individual organisms are best thought of as adaptation-executers rather than as fitness-maximizers.'' I put it at the top and use it as my title.

Fifty thousand years ago, taste buds led people to sugar and fat, which were scarce. Today calories are cheap in rich countries, and the same taste buds lead people to ice cream. No one who deliberately set out to maximize their genetic fitness, I say, would eat a cookie unless they were starving. \nb{Biologists call a case like this an evolutionary mismatch: a trait that was advantageous becomes harmful when the environment changes. The post does not use the term.}

To explain, I take a screwdriver. Its cause is the designer's mind. Its shape is the only one of these things found in the tool. Its consequence is what it actually does, which may be cutting the user's hand. Its meaning is what a user takes it to be for, which may be a weapon or a chisel. The tool knows nothing of its purpose, and will not reshape itself to fit a flat-head screw.

Then I apply the same four headings to the taste bud. Its cause is a long line of ancestors who liked sugar and had children, ``2763 generations'' of them until the gene spread. \nb{The figure matches the calculation in Yudkowsky's ``Evolutions Are Stupid (But Work Anyway)'' for a gene with a 1\% advantage in a population of a million.} Its shape is a sensor linked to reinforcement. Its consequence, today, can run from chocolate to getting fat to having fewer children. Its meaning is up to you; I like chocolate and wish it were less harmful. \nb{Separating the history that explains a trait from the mechanism that makes it work is what biologists call ultimate and proximate causation, terms Ernst Mayr popularized in 1961. The post does not use them.}

It is not quite right, I add, to say that people do today what would have spread their genes among hunter-gatherers, because no one asks that question, and chocolate has no ancestral counterpart. ``Therefore it is said,'' I conclude, and repeat the epigraph.
